% Sections 10.1 to 10.5, translated from sv/chapters/10/theory.tex.

\section{Exponential functions and equations}
\label{c10:sec:exponential}
An equation with an \textbf{exponential function} $f(x) = a^x$ ($a > 0$, $a \neq 1$) can be solved
for $x$ with logarithms.

\textbf{Solving $a^x = b$:} take the logarithm of both sides.
\[
  x \cdot \log a = \log b \quad \Rightarrow \quad x = \frac{\log b}{\log a}
\]

\section{Logarithms}
\label{c10:sec:logaritmer}
The \textbf{base-10 logarithm} $\log_{10}$, usually written $\log$, is defined by
\[
  \log_{10}(10^x) = x \quad \Leftrightarrow \quad 10^{\log x} = x.
\]
The \textbf{natural logarithm} $\ln$ has the base $e \approx 2.718$:
\[
  \ln(e^x) = x \quad \Leftrightarrow \quad e^{\ln x} = x
\]

\section{The laws of logarithms}
\label{c10:sec:logaritmlagar}
Logarithms obey the following laws:
\begin{gather*}
  \log(ab) = \log a + \log b \\
  \log\left(\frac{a}{b}\right) = \log a - \log b \\
  \log(a^n) = n \cdot \log a
\end{gather*}
\textbf{Change of base:}
\[
  \log_a b = \frac{\log b}{\log a} = \frac{\ln b}{\ln a}
\]

\section{Decibels (dB)}
\label{c10:sec:decibel}
Decibels are used in electrical engineering to give gain and attenuation on a logarithmic scale.

\textbf{Voltage gain:}
\[
  G_{\text{dB}} = 20 \log_{10}\left(\frac{U_{\text{out}}}{U_{\text{in}}}\right)
\]
Conversely:
\[
  \frac{U_{\text{out}}}{U_{\text{in}}} = 10^{G_{\text{dB}}/20}
\]
\textbf{Power gain:}
\[
  G_{\text{dB}} = 10 \log_{10}\left(\frac{P_{\text{out}}}{P_{\text{in}}}\right)
\]
\textbf{Level in dBV}, decibels relative to $1\,\text{V}$:
\[
  U_{\text{dBV}} = 20 \log_{10}\left(\frac{U_{\text{RMS}}}{1\,\text{V}}\right)
\]

\needspace{12\baselineskip}
\section{Worked examples}
\label{c10:sec:typexempel}

\begin{exempel}[Solving logarithmic equations]
Solve \textbf{a)} $3^x = 81$, \textbf{b)} $2^{x-1} = 64$ and \textbf{c)} $e^{x-2} = 150$.

\textbf{a)} $3^x = 3^4 \Rightarrow x = 4$. Alternatively, $x = \frac{\log 81}{\log 3} = 4$.

\textbf{b)}
\[
  (x-1)\log 2 = \log 64 \quad \Rightarrow \quad x - 1 = \frac{\log 64}{\log 2} = 6
  \quad \Rightarrow \quad x = 7
\]

\textbf{c)}
\[
  x - 2 = \ln 150 \approx 5.01 \quad \Rightarrow \quad x \approx 7.01
\]
\end{exempel}

\begin{exempel}[Half-life of a battery charge]
A battery loses half its charge in $30$ hours:
\[
  u(30) = U_0 \cdot a^{30} = 0.5\,U_0
\]
\textbf{Calculate the growth factor $a$:}
\[
  a^{30} = 0.5 \quad \Rightarrow \quad a = 0.5^{1/30} \approx 0.977
\]
\textbf{When is $20\,\%$ left?}
\[
  a^t = 0.2 \quad \Rightarrow \quad t = \frac{\log 0.2}{\log a}
  = \frac{\log 0.2}{\log 0.5^{1/30}} \approx 69.7\,\text{h}
\]
Calculate with the unrounded value of $a$; with $0.977$ the answer becomes $69.2\,\text{h}$.
\end{exempel}

\begin{exempel}[Linear voltage gain]
Two signals have the levels $L_1 = 20\,\text{dB}$ and $L_2 = 46\,\text{dB}$.
\begin{gather*}
  G_{\text{dB}} = L_2 - L_1 = 26\,\text{dB} \\
  G_{\text{lin}} = 10^{26/20} = 10^{1.3} \approx 20
\end{gather*}
\end{exempel}

\begin{exempel}[dBV to volts]
A sinusoidal voltage has the level $31.0\,\text{dBV}$.
\begin{gather*}
  U_{\text{RMS}} = 10^{31.0/20} \approx 35.5\,\text{V} \\
  \abs{U} = U_{\text{RMS}} \cdot \sqrt{2} \approx 35.5 \cdot 1.414 \approx 50.2\,\text{V}
\end{gather*}
\end{exempel}
